% section_2.tex — Approach (questionnaire design + streaming platform)

\section{Approach}
\label{sec:approach}

We start from existing third-party renderer implementations for the
three NVS methods we evaluate --- \texttt{pyngp} for
i-NGP~\citep{mueller2022instantngp},
\texttt{diff-nexel-rasterization} for Nexels~\citep{nexels}, and
\texttt{gsplat}~\citep{ye2024gsplat} for 3D Gaussian
Splatting~\citep{kerbl20233dgs}. What those projects do not ship is the
instrument: there is no agreed-upon questionnaire for asking a
participant how good a real-time 3D rendering feels, and no
off-the-shelf platform that runs such a questionnaire while the
participant interactively views a live render. The contribution of this
work is, in order of importance, the design of that questionnaire
(\S\ref{sec:questionnaire}) and the streaming harness that lets it run
(\S\ref{sec:platform}).

\subsection{Questionnaire design}
\label{sec:questionnaire}

\paragraph{Decomposing ``quality''.}
A single overall ``how good is this rendering'' rating is hard to
interpret across scenes and methods: it averages over different failure
modes, and the participant's internal weighting is opaque. We instead
decompose perceived quality into three sub-aspects that we elicit
separately:
\begin{enumerate}
    \item \emph{Foreground detail} --- sharpness and fidelity on the
    nominal subject of the captured photographs.
    \item \emph{Background detail} --- sharpness, fidelity, and absence
    of floater artifacts on the surrounding regions, which NVS methods
    tend to struggle with more than the foreground.
    \item \emph{Motion smoothness} --- continuity of the rendering during
    free navigation. This is the one aspect that \emph{requires}
    interactive viewing to evaluate at all.
\end{enumerate}
The first two are familiar from still-image quality work; the third is
the explicit reason this study exists as a streaming harness rather than
an image grader. The three axes also align with the kinds of design
tradeoffs an NVS researcher actually reasons about --- model capacity
allocated to subject vs.\ surroundings, and the speed cost of additional
parameters or splats.

\paragraph{Sliders rather than Likert.}
Each aspect is rated on a continuous 0--100 slider with the endpoints
labeled (``poor'' and ``excellent'') and the midpoint unlabeled. We
chose continuous sliders over a 5-point Likert scale for two reasons:
(i) sliders avoid the anchor-driven clustering that compresses Likert
distributions toward one or two modes, and (ii) on a 0--100 scale a
single within-participant slider can be correlated directly against an
offline metric on the same axis, which a 5-bucket categorical cannot.

\paragraph{Two-method comparison: independent viewports, forced A/B,
aspect citation.}
For the two-method comparison the participant sees two
\emph{independently orbitable} viewports side by side, with method
identities hidden behind anonymous badges (only the \texttt{test}
initials hook reveals them, and is never used in real sessions).
Independence is deliberate: a single baked camera trajectory would
dictate which artifacts the participant ever sees, but pilot sessions
confirmed that participants want to \emph{probe} the scene --- zooming
on suspicious regions, viewing the subject from oblique angles --- and
the artifacts they probe (popping, smoothness, floaters at grazing
angles) are precisely the ones a fixed trajectory would either hide or
showcase. After exploring a pair, the participant fills out per-side
foreground / background / smoothness sliders, then makes a forced A/B
overall preference (``A'', ``B'', or explicit tie), then selects from a
multi-select list \emph{which} aspect(s) shaped that choice. The
slider-then-A/B redundancy is intentional: we collect both a graded
within-method rating and a categorical between-method judgment per
trial, and the aspect multi-select converts an aggregate preference
into something mechanistic --- which aspect, on which scene, swung the
choice.

\paragraph{Single-method convergence: ratings in isolation, reference
anchor.}
The convergence case uses the same primitives in a different shape. The
participant sees one viewport, rendering one stimulus (here: a 3DGS
checkpoint at a given training-iteration count) at a time; the active
stimulus is swapped without a process restart, so the only thing that
changes between consecutive stimuli is the model state. We deliberately
rate each stimulus in isolation rather than as A/B against the previous
one, since the stimuli are not always commensurable on a single axis of
merit and a forced choice would discard graded information. Per
stimulus we collect two sliders --- overall quality and motion
smoothness --- on the same 0--100 scale as the comparison case. To
give the participant a ground-truth anchor for what the scene
\emph{should} look like, we play a low-frame-rate slideshow of the
training photographs for the scene once before its rating block; an
NVS rendering on its own provides no reference for ``how detailed
should this actually be''.

\subsection{Streaming platform}
\label{sec:platform}

The streaming pipeline exists to let either questionnaire above run
while the participant views a live render rather than a pre-baked
trajectory, despite participants connecting from their own
(typically GPU-less) laptops, and despite the underlying renderers
having different native rendering stacks --- the two constraints
flagged in \S\ref{sec:background}. The platform is, by intent, the
smaller half of the contribution: it is what allowed us to ship the
questionnaire, not a research artifact in itself.

\paragraph{Pipeline.}
Each renderer is wrapped in a thin process
(\texttt{interactive\_renderer.py},
\texttt{nexels\_renderer.py}, \texttt{gsplat\_renderer.py}) that
consumes camera-pose deltas on a per-method WebSocket and pushes H.264
over RTSP to a \texttt{mediamtx} instance, which relays to the browser
over WebRTC (WHEP). A single browser-facing WebSocket multiplexer
(\texttt{ws\_mux.py}) routes pose-delta messages to the right renderer
based on URL path, and a static HTTP server (\texttt{serve.py}) hosts
the viewer pages and accepts rating submissions. All server-side
processes colocate on a single AWS L4 GPU instance.
Figure~\ref{fig:architecture} shows the resulting data flow.

\begin{figure}[t]
  \centering
  \includegraphics[width=\linewidth]{architecture}
  \caption{Streaming pipeline. The three renderers conform to a common
  contract (RTSP push for frames out, per-method WebSocket for events
  in); \texttt{ws\_mux.py} and \texttt{mediamtx} between them keep the
  browser renderer-agnostic.}
  \label{fig:architecture}
\end{figure}

\paragraph{Plug-and-play seam.}
The principal design choice on the platform side is the contract every
renderer must satisfy: (a) accept camera-pose deltas as JSON messages
on a per-method WebSocket port; (b) push H.264 over RTSP at the
configured target resolution and frame rate (1440$\times$810@30 for the
comparison case, 1920$\times$1080@60 for the convergence case); and
(c) optionally accept a \texttt{swap\_state} message used by the
convergence case to flip the active checkpoint without restarting the
renderer process. Any process that meets this contract drops in by
registering a route in \texttt{ws\_mux.py} and a corresponding RTSP
path in \texttt{mediamtx.yml}. The browser never has to know which
underlying method is active --- the same viewer HTML works for any
renderer that honors the contract, and the marginal cost of adding a
third method is a few lines of configuration rather than a viewer
rewrite.

\paragraph{What did not work.}
Two earlier designs are preserved in \texttt{archive/}. A fully
pre-rendered evaluation built on canned camera trajectories rendered to
MP4 surfaced the exact limitation that motivated the streaming
approach --- participants could not probe the artifacts they wanted to
talk about. A second attempt used Xvfb plus each renderer's native UI,
with \texttt{event\_relay.py} translating browser events into
\texttt{xdotool} keystrokes; this worked end-to-end but the
screen-capture path added tens of milliseconds, and the design did not
generalize to running two renderers concurrently because each demanded
its own X display. The current architecture replaced both.
